\begin{abstract}
Although large-scale unsupervised language models (LMs) acquire extensive world knowledge and develop some reasoning abilities, their behavior is challenging to control precisely because their training is entirely unsupervised. Current approaches to enhance controllability involve gathering human judgments about the relative quality of different model outputs and adapting the LM using this preference data, commonly through reinforcement learning from human feedback (RLHF). Nonetheless, RLHF presents difficulties, as it typically requires first training a reward model to mimic human preferences, then applying reinforcement learning to fine-tune the unsupervised LM to optimize this proxy reward, all while preventing excessive deviation from the pretrained model. This work introduces a novel reward model parameterization within the RLHF framework, which allows for direct extraction of the corresponding optimal policy in closed form. As a result, we can address the conventional RLHF objective via a straightforward classification loss. The new approach, termed \textit{Direct Preference Optimization} (DPO), offers stable and strong performance and is computationally efficient, doing away with the necessity to sample from the LM during fine-tuning and greatly reducing the burden of hyperparameter tuning. Empirical results demonstrate that DPO fine-tunes LMs to follow human preferences on par with or better than established techniques. Importantly, DPO surpasses PPO-based RLHF in steering the sentiment of generated text and performs comparably or better for summarization and single-turn conversation tasks, while also being far simpler to implement and train.
\end{abstract}

\section{Introduction}
Large-scale unsupervised language models (LMs), when trained on massive corpora, often display unexpected and advanced capabilities~\citep{chowdhery2022palm, brown2020language, touvron2023llama,bubeck2023sparks}. Nevertheless, the training data for these models is composed of human-generated content that reflects a diverse array of intentions, competencies, and objectives. Some of these objectives or skills may be inappropriate or suboptimal for models to replicate; for example, although it is advantageous for an AI programming assistant to \textit{recognize} frequent programming errors so that it can correct them, we would prefer the model to \emph{favor} the (sometimes infrequent) examples of proficient coding present in the data when actually generating code. In a similar vein, while it is useful for a language model to \textit{possess knowledge of} widespread misconceptions—such as one held by 50\% of individuals—it is highly undesirable for the model to repeat the misconception as fact in 50\% of its relevant outputs. Thus, effectively guiding a model to exhibit the \emph{preferred responses and behaviors}, as distinct from its broader \textit{knowledge base and capabilities}, is a key component in designing AI systems that are reliable, effective, and manageable \citep{ouyang2022training}. Although common practice involves steering LMs toward alignment with human preferences via reinforcement learning (RL), we demonstrate that the RL-oriented objective underlying such methods can in fact be directly optimized using a straightforward binary cross-entropy loss, thereby streamlining the process of preference learning significantly.

Broadly, the prevailing paradigm for imbuing language models with desired characteristics involves utilizing carefully constructed datasets that encode human preferences for safe and useful behavior. This preference optimization generally takes place after a preliminary phase of extensive unsupervised pre-training on large-scale text datasets. The simplest method for this task—supervised fine-tuning on human-curated demonstrations of ideal outputs—is typically surpassed by a family of approaches centered around reinforcement learning from human (or AI) feedback (RLHF/RLAIF; \citep{christiano2017deep,bai2022constitutional}). In RLHF, a reward model is first trained on data reflecting human preferences, after which the policy of the language model is optimized using RL so as to produce outputs that maximize the reward, while constraining the policy to remain close to the pre-trained model. Despite the effectiveness of RLHF in producing models with state-of-the-art conversational and programming proficiency, the RLHF framework is considerably more intricate than supervised fine-tuning. It requires training multiple language models and iterative sampling from the policy during training, which results in substantial computational overhead.

This work presents a method for directly training language models to align with human preferences without the need for explicit reward models or reinforcement learning techniques. We introduce \textit{{Direct Preference Optimization} (DPO)}, a procedure that effectively targets the same objective as prevalent RLHF strategies (specifically, maximizing reward subject to a KL-divergence penalty), yet remains simple to deploy and efficient to train. Conceptually, DPO operates by increasing the log probability ratio in favor of preferred responses over non-preferred ones, but crucially incorporates adaptive, instance-specific weighting to counteract the over-optimization and degradation issues we observed with straightforward probability ratio objectives. Similar to established algorithms, DPO draws upon a theoretical framework for preferences (for example, the Bradley-Terry model; \cite{bradley1952rankanalysis}) to evaluate how well a reward function matches observed preference judgments. The key difference is that, whereas existing pipelines utilize this preference model to form a loss for training a reward function (which then guides policy optimization), DPO applies a reparameterization to express the preference loss directly in terms of the policy. Consequently, using a corpus of human-annotated preference comparisons, DPO optimizes a policy through a binary cross entropy loss, resulting in the best-fit policy for an implicit reward aligned with those preferences.

The principal contribution of this paper is the introduction of Direct Preference Optimization (DPO), which is a reinforcement learning-free framework for preference-based language model fine-tuning. Our empirical studies indicate that DPO achieves performance on par with leading alternatives, including PPO-based RLHF, across a range of preference-driven benchmarks such as sentiment control, text summarization, and dialogue modeling, and scales to models with up to 6B parameters.

Language models trained in a self-supervised manner at increasingly large scales are capable of performing certain tasks in a zero-shot setting \citep{radford2019language} or with minimal demonstration via few-shot prompting \citep{gpt3,megatron,chowdhery2022palm}. Nonetheless, their capabilities on downstream applications and alignment with human intent are often greatly enhanced through a process of fine-tuning on curated datasets consisting of instructions paired with human-generated completions \citep{mishra-etal-2022-cross,sanh2022multitask,chung2022scaling,thoppilan2022lamda}. This strategy, known as `instruction-tuning', allows large language models (LLMs) to extrapolate to novel instructions not explicitly present during tuning, thereby broadening their practical utility \citep{chung2022scaling}. While instruction-tuning has achieved notable advances, it is frequently more feasible to gather \textit{relative} assessments of output quality from non-experts than it is to obtain authoritative demonstrations. In response, later studies have focused on adapting LLMs to human preference datasets, yielding gains in diverse areas including translation \citep{kreutzer-etal-2018-reliability}, summarization \citep{stiennon2022learning,ziegler2020finetuning}, narrative generation \citep{ziegler2020finetuning}, and instruction adherence \citep{ouyang2022training,ramamurthy2023is}. These approaches typically begin by training a neural reward model that aligns with the preference data, often employing a framework such as the Bradley-Terry model \citep{bradley1952rankanalysis} to represent comparisons. The language model is then updated to optimize for this learned reward using reinforcement learning algorithms, such as REINFORCE \citep{williams1992reinforce}, proximal policy optimization (PPO; \cite{schulman2017proximal}), or similar variants \citep{ramamurthy2023is}. A related area of research harnesses LLMs that have been fine-tuned via instruction-following with human feedback, employing them to create additional synthetic datasets focused on attributes like safety or harmlessness \citep{bai2022constitutional}, utilizing only broad textual rubrics as weak human supervision for their labeling processes. Such developments blend advances from two domains: one focusing on reinforcement learning as a means of training language models for multiple objectives~\citep{Ranzato2015SequenceLT,paulus2018a,wu2018learning}, and another dedicated to methodologies for learning from human preference data \citep{christiano2017deep,kupcsik2018learning}. Despite the intuitive advantages of optimizing with relative human preferences, applying reinforcement learning to fine-tune large language models remains logistically challenging; the present work introduces a theoretically-grounded alternative for optimizing over relative preferences without the use of RL.

Outside of language-related scenarios, the study of learning policies from preferences has been extensively explored within both the bandit and reinforcement learning paradigms, yielding a variety of methodologies. When contextual bandit algorithms use preferences or rankings over actions instead of explicit reward feedback, the problem is referred to as a contextual dueling bandit (CDB; \cite{yue2012karmed,dudik2015contextual}). In these contexts where absolute rewards are unavailable, the theoretical foundation of CDBs relies on the concept of a \textit{von Neumann winner}: a policy whose expected win probability is no less than 50\% when compared against any other policy \citep{dudik2015contextual}. Notably, CDB frameworks receive preference annotations in an online manner, whereas in settings involving human preference learning, models typically leverage a fixed offline dataset consisting of preference-labeled action pairs \citep{yan2022human}. Along similar lines, \textit{preference-based RL} (PbRL) learns by observing binary preferences dictated by an \textit{unknown} `scoring' or utility function rather than direct rewards \citep{BusaFekete2014,ruiz2023dueling}. Multiple algorithmic solutions exist for PbRL—including those that effectively utilize off-policy preference data—but most require the explicit construction of a latent reward or scoring function prior to optimization \citep{jain2013learning,BusaFekete2014,christiano2017deep,sadigh2017active,kupcsik2018learning}. Contrasting with these approaches, we introduce a unified method that directly learns a policy to align with observed preferences, eliminating the need for separate reward model estimation and policy optimization steps.

\section{Preliminaries}\label{section:prelims}

We begin by summarizing the RLHF framework as presented by \citeauthor{ziegler2020finetuning} and further developed in works such as \citep{stiennon2022learning, bai2022training, ouyang2022training}. This procedure typically unfolds in three primary steps: (1) supervised fine-tuning (SFT), (2) collection of preference data coupled with reward model estimation, and (3) reinforcement learning-based optimization.

\textbf{SFT}: The RLHF pipeline generally commences with supervised adaptation of a pre-trained language model (LM) using high-quality, task-specific data (e.g., for dialogue or summarization), thereby yielding an initial model denoted $\pi^\text{SFT}$.

\textbf{Reward Modelling Phase}: Next, the fine-tuned SFT model is conditioned on input prompts $x$ to generate candidate response pairs $(y_1, y_2)\sim \pi^\text{SFT}(y \mid x)$. These pairs are then presented to human raters, who indicate a preference between the two responses—formally denoted as $y_w\succ y_l \mid x$, where $y_w$ and $y_l$ refer to the selected (preferred) and non-selected (dispreferred) outputs, respectively, for each prompt. The underlying assumption is that such preferences are determined by some unobservable reward function $r^*(y, x)$, which is not directly accessible. Various preference modeling approaches exist, with the Bradley-Terry (BT) \cite{bradley1952rankanalysis} model frequently adopted (though the more general Plackett-Luce models \citep{plackett1975analysis, luce2012individual} can also be integrated when preferences over ranked lists are available). Within the BT framework, the likelihood of the human preference distribution $p^*$ can be represented as follows:

\begin{equation}\label{eq:bradley-terry}
    p^*(y_1\succ y_2 \mid x)=\frac{\exp\left(r^*(x, y_1)\right)}{\exp\left(r^*(x, y_1)\right) + \exp\left(r^*(x, y_2)\right)}.
\end{equation}

Given a fixed dataset of comparison samples $\mathcal{D}=\bigl\{x^{(i)}, y_w^{(i)}, y_l^{(i)}\bigr\}_{i=1}^N$ drawn from $p^*$, we can define a parameterized reward model $r_{\phi}(x, y)$ and fit its parameters through maximum likelihood estimation. Interpreting the task as a binary classification, we arrive at the negative log-likelihood loss:

\begin{equation}\label{eq:reward_model}
    \mathcal{L}_R(r_{\phi}, \mathcal{D}) = -\mathbb{E}_{(x, y_w, y_l)\sim \mathcal{D}}\bigl[\log \sigma(r_{\phi}(x, y_w)- r_{\phi}(x, y_l))\bigr]
\end{equation}

where $\sigma$ denotes the sigmoid function. For large language models, the function $r_{\phi}(x, y)$ is typically initialized from the SFT model $\pi^\text{SFT}(y \mid x)$, augmented by appending a linear layer after the last transformer block that outputs a scalar reward estimate \cite{ziegler2020finetuning}. To stabilize the reward model, prior work often normalizes its outputs so that $\mathbb{E}_{x,y\sim \mathcal{D}}\left[r_\phi(x, y)\right] = 0$ for each $x$.

\textbf{RL Fine-Tuning Phase}: During the reinforcement learning stage, the trained reward model acts as a feedback mechanism for the language model policy. In accordance with previous research~\citep{jaques2017sequence, jaques2020human}, the objective to be maximized is expressed as

\begin{equation}\label{eq:RL}
\max_{\pi_{\theta}}  \mathbb{E}_{x\sim \mathcal{D}, y\sim \pi_{\theta}(y \mid x)}\bigl[r_{\phi}(x, y)\bigr] - \beta\mathbb{D}_{\textrm{KL}}\bigl[\pi_{\theta}(y\mid x)\mid \mid \pi_\text{ref}(y\mid x)\bigr],
\end{equation}

where $\beta$ modulates the penalty for diverging from the reference policy $\pi_\text{ref}$, which is typically the SFT-initialized model $\pi^\text{SFT}$. The policy $\pi_\theta$ is also commonly initialized from $\pi^\text{SFT}$. The KL penalty is critical as it ensures that the policy does not stray excessively from distributions where the reward model remains accurate, while also maintaining sample diversity and mitigating mode-collapse towards singular, high-reward outputs. Owing to the discrete output space of language generation, this objective cannot be optimized via standard gradient methods and instead necessitates reinforcement learning. Commonly, the reward function is defined as ${r(x, y) = r_{\phi}(x, y) -\beta (\log \pi_{\theta}(y\mid x) - \log \pi_\text{ref}(y\mid x))}$, and optimization proceeds using PPO \cite{schulman2017proximal}, as seen in \cite{ziegler2020finetuning, stiennon2022learning, bai2022training, ouyang2022training}.

\section{Direct Preference Optimization}\label{sec:DPO}

Recognizing the computational and algorithmic difficulties involved in applying traditional reinforcement learning techniques to large language models, our aim is to establish a streamlined method for optimizing policies directly from preference data. In contrast to established RLHF workflows—which require first learning an explicit reward function before subsequent RL-based optimization—we exploit a specific reward model formulation that admits a closed-form solution for its optimal policy, thereby circumventing the need for reinforcement learning iterations. Our central insight is to utilize a direct analytical relationship between reward models and the corresponding optimal policies, facilitating a transformation of reward-based loss functions into policy-based objectives. This reparameterization obviates the need for a separately-trained reward model, yet preserves alignment with standard preference-based frameworks like the Bradley-Terry model. Effectively, the policy network simultaneously models the language policy and an implicit reward function.

\textbf{Deriving the DPO objective.} To begin, consider the standard reinforcement learning objective (Eq.~\ref{eq:RL}) with a general reward function $r$, as commonly adopted in earlier research. Building on established results~\citep{peters2007reinforcement, peng2019advantage, korbak2022reinforcement, go2023aligning}, it can be readily shown that the optimal solution to the KL-regularized reward maximization problem in Eq.~\ref{eq:RL} is given by:

\begin{equation}\label{eq:op_policy}
    \pi_r(y\mid x) = \frac{1}{Z(x)}\pi_\text{ref}(y\mid x)\exp\left(\frac{1}{\beta}r(x, y)\right),
\end{equation}

with the normalization constant $Z(x) =\sum_{y}\pi_\text{ref}(y\mid x)\exp\left(\frac{1}{\beta}r(x, y)\right)$. Full details are provided in Appendix \ref{app:derivation1}. Notably, even with an MLE-based approximation $r_{\phi}$ to the true reward $r^*$, direct computation of $Z(x)$ can be computationally intensive \citep{korbak2022reinforcement, go2023aligning}, limiting the practicality of this parameterization. To address this, Eq.~\ref{eq:op_policy} can be manipulated to reformulate the reward in terms of the optimal policy $\pi_r$, the reference policy $\pi_\text{ref}$, and the (unknown) $Z(x)$. Taking the logarithm of Eq.~\ref{eq:op_policy} and rearranging terms yields:

\begin{equation}\label{eq:main_eq}
    r(x,y) =\beta \log \frac{\pi_r(y\mid x)}{\pi_\text{ref}(y\mid x)} + \beta \log Z(x).
\end{equation}

This form can be applied to both the ground-truth reward $r^*$ and its associated optimal policy $\pi^*$. Importantly, the Bradley-Terry preference model only depends on differences in reward between pairs of outputs; that is, ${p^*(y_1 \succ y_2 \mid x) = \sigma(r^*(x, y_1) - r^*(x, y_2))}$. Substituting the new expression from Eq.~\ref{eq:main_eq} for $r^*(x,y)$ into the preference model Eq.~\ref{eq:bradley-terry} reveals that $Z(x)$ cancels out, enabling us to write the probability of human preference solely in terms of the optimal policy $\pi^*$ and the reference $\pi_\text{ref}$. Therefore, under the Bradley-Terry model, the optimal RLHF policy $\pi^*$ must satisfy:

\begin{equation}\label{eq:objective}
    p^*(y_1\succ y_2 \mid x)=\frac{1}{1 + \exp\left(\beta \log \frac{\pi^*(y_2\mid x)}{\pi_\text{ref}(y_2\mid x)} - \beta \log \frac{\pi^*(y_1\mid x)}{\pi_\text{ref}(y_1\mid x)}\right)}
\end{equation}

See Appendix~\ref{app:derivation2} for details. Although the formulation above is based on the Bradley-Terry model, similar derivations apply for the more general Plackett-Luce models~\citep{plackett1975analysis, luce2012individual}; see Appendix~\ref{app:plackett_luce_models}.

This reparameterization permits the use of the optimal policy, instead of the reward function, to model human preference probabilities. Consequently, we can directly optimize a likelihood objective over the policy parameters $\pi_\theta$. Mirroring the construction in reward modeling (cf. Eq.~\ref{eq:reward_model}), we arrive at the following policy learning objective:

\begin{equation}\label{eq:optimum_model}
    \mathcal{L}_\text{DPO}(\pi_{\theta}; \pi_\text{ref}) = -\mathbb{E}_{(x, y_w, y_l)\sim \mathcal{D}}\left[\log \sigma \left(\beta \log \frac{\pi_{\theta}(y_w\mid x)}{\pi_\text{ref}(y_w\mid x)} - \beta \log \frac{\pi_{\theta}(y_l\mid

In this formulation, we model an implicit reward through an alternative parameterization in which the optimal policy corresponds directly to $\pi_\theta$. Additionally, since this approach amounts to learning a reparameterized Bradley-Terry model, it inherits favorable theoretical characteristics, including consistency under appropriate assumptions regarding the distribution of preference data \cite{bong2022generalized}. Section~\ref{sec:theory} elaborates further on the theoretical guarantees of DPO and compares them with existing literature.

\textbf{How does the DPO update function?} To elucidate the workings of DPO, it is instructive to examine the gradient of its loss function, $\mathcal{L}_\text{DPO}$. The gradient with respect to model parameters $\theta$ is expressed as follows:

\begin{multline*}\label{eq:gradient}
    \nabla_\theta \mathcal{L}_\text{DPO}(\pi_\theta;\pi_\text{ref}) = \\ -\beta\mathbb{E}_{(x, y_w, y_l) \sim \mathcal{D}} \bigg[\underbrace{\sigma(\hat{r}_\theta(x, y_l) - \hat{r}_\theta (x, y_w))}_\text{emphasizes examples where reward ranking is incorrect}\bigg[\underbrace{\nabla_\theta\log \pi(y_w \mid x)}_\text{promote $y_w$ occurrence} - \underbrace{\nabla_\theta\log\pi(y_l \mid x)}_\text{reduce $y_l$ occurrence}\bigg]\bigg],
\end{multline*}

Here, $\hat{r}_\theta(x, y) = \beta \log \frac{\pi_\theta(y \mid x)}{\pi_\text{ref}(y \mid x)}$ denotes the reward signal implicitly determined by the candidate policy $\pi_\theta$ relative to the reference policy $\pi_\text{ref}$ (see also Section~\ref{sec:theory}). Conceptually, taking the gradient of $\mathcal{L}_\text{DPO}$ has the effect of enhancing the likelihood of the preferred outputs $y_w$ while diminishing the probability of the less desirable completions $y_l$. Notably, each training instance is weighted in proportion to the extent that the implicit reward model $\hat{r}_\theta$ overvalues the less preferred answer, modulated by the parameter $\beta$, thus quantifying the reward model’s misordering while incorporating the influence of the KL regularization. Empirical findings underscore the necessity of this weighting mechanism: omitting the weighting factor results in substantial degradation of model performance, as illustrated in Appendix Table~\ref{tab:unlikelihood_generations}.

\textbf{DPO overview.} The typical workflow for DPO consists of: 1) Drawing sample completions $y_1, y_2 \sim \pi_\text{ref}(\cdot \mid x)$ for each prompt $x$, then annotating them with human preference labels to form an offline preference dataset $\mathcal{D} = \{x^{(i)}, y_w^{(i)}, y_l^{(i)}\}_{i=1}^N$; and 2) training the language model $\pi_\theta$ by minimizing the loss $\mathcal{L}_\text{DPO}$, using the specified $\pi_\text{ref}$, dataset $\mathcal{D}$, and a chosen value for $\beta$.

In many applications, leveraging existing, publicly available preference datasets is preferable to creating new samples and collecting human annotations. Since these datasets are typically derived using $\pi^\text{SFT}$, we set $\pi_\text{ref} = \pi^\text{SFT}$ when such a model is accessible. If $\pi^\text{SFT}$ is absent, we construct $\pi_\text{ref}$ by maximizing the likelihood over the preferred responses ${(x, y_w)}$, meaning ${\pi_\text{ref} = \argmax_{\pi}\mathbb{E}_{x, y_w \sim \mathcal{D}}\left[\log \pi(y_w \mid x)\right]}$. Adopting this approach helps alleviate the mismatch between the inaccessible true reference distribution and the practical $\pi_\text{ref}$ employed in DPO. Additional specifics on implementation choices and hyperparameter settings are available in Appendix~\ref{app:implementation}.

\section{Theoretical Analysis of DPO} In this section, we offer a deeper conceptual understanding of DPO, provide formal justification, and explain how DPO addresses limitations found in RLHF actor-critic frameworks (for instance, PPO~\cite{schulman2017proximal}).

\label{sec:theory}

\subsection{Your Language Model Is Secretly a Reward Model} DPO sidesteps the necessity of explicitly learning a reward function and carrying out separate RL policy optimization, instead employing a unified maximum likelihood objective. Notably, the objective in Eq. \ref{eq:main_eq} matches the form of a Bradley-Terry model, where the reward is parameterized as $r^*(x, y) = \beta \log\frac{\pi^*_\theta(y \mid x)}{\pi_\text{ref}(y \mid x)}$, and the parametric model $\pi_{\theta}$ is directly optimized—this corresponds to reward model fitting under a variable transformation as described in Eq. \ref{eq:reward_model}. This section formalizes the theoretical basis for this reparameterization, demonstrates that it imposes no restrictions on the set of expressible reward models, and ensures that the optimal policy can be exactly recovered. We begin by introducing an equivalence relation over reward functions.

\begin{definition}
Two reward functions $r(x, y)$ and $r'(x, y)$ are said to be equivalent if and only if there exists some function $f$ such that ${r(x, y) - r'(x, y) = f(x)}$.
\end{definition}

It is straightforward to verify that this relationship constitutes an equivalence relation, which divides the set of all reward functions into distinct equivalence classes. With this in mind, the following lemmas can be established:

\begin{lemma}\label{lemma:same_prefrence}
Within the Plackett-Luce framework—including, in particular, the Bradley-Terry model—any two reward functions from the same equivalence class produce identical preference distributions.
\end{lemma}

\begin{lemma}\label{lemma:same_policy}
    Any pair of reward functions belonging to the same equivalence class will lead to the same optimal policy when solving the constrained reinforcement learning (RL) problem.
\end{lemma}

Since the proofs for these statements are elementary, we postpone them to Appendix \ref{app:lemma1}. The first lemma is a commonly acknowledged issue regarding the non-uniqueness (under-specification) present in the Plackett-Luce family of models \cite{plackett1975analysis}. This non-identifiability necessitates the use of additional constraints to ensure meaningful maximum likelihood estimation, as referenced in Eq. \ref{eq:reward_model} \cite{bong2022generalized}. The second lemma guarantees that every reward function within a given equivalence class results in the same optimal policy. As such, for purposes of optimization, it suffices to recover any reward function from the optimal equivalence class. The main result, which we demonstrate in Appendix~\ref{app:thm1}, is given below:

\begin{theorem}\label{thm:main}
    Provided some mild conditions hold, each reward equivalence class consistent with the Plackett-Luce (and specifically the Bradley-Terry) model admits a reparameterization of the form ${r(x, y) = \beta \log \frac{\pi(y\mid x)}{\pi_\text{ref}(y\mid x)}}$, where $\pi(y\mid x)$ is some model and $\pi_\text{ref}(y \mid x)$ is a specified reference distribution.
\end{theorem}

\begin{sproof}
    Let $r(x, y)$ denote any reward function that defines a corresponding optimal policy $\pi_r(y \mid x)$, given by Eq. \ref{eq:op_policy}. We aim to show that one can construct a representative within the equivalence class of $r$ using the stated reparameterization. To do this, define the following transformation:
\begin{equation}
    f(r; \pi_\text{ref}, \beta)(x, y) = r(x, y) - \beta\log\sum_{y}\pi_\text{ref}(y\mid x)\exp\left(\frac{1}{\beta}r(x, y)\right)
\end{equation}
This operator $f$ effectively rescales the reward by subtracting the logarithm of the partition function associated with $\pi_r$. Since the normalization term depends solely on $x$, $f(r; \pi_\text{ref}, \beta)(x, y)$ remains within the same equivalence class as $r(x, y)$. Substituting $r$ using the right-hand side of Eq.~\ref{eq:main_eq} (valid for any $r$) yields $f(r; \pi_\text{ref}, \beta)(x, y) = \beta \log \frac{\pi_r(y\mid x)}{\pi_\text{ref}(y\mid x)}$. This demonstrates that $f$ constructs a representative in the equivalence class of $r$ with the desired structure, thereby showing that the proposed reparameterization covers all relevant cases in the reward model.
\end{sproof}

An alternative perspective on Theorem~\ref{thm:main} is that it identifies the precise reward function chosen by the DPO reparameterization within each equivalence class—specifically, the reward function that satisfies:

\begin{equation}\label{eq:lag_p}
     \sum_{y}\underbrace{\pi_\text{ref}(y\mid x)\exp\left(\frac{1}{\beta}r(x, y)\right)}_{=\pi(y\mid x)\text{, using Thm.~\ref{thm:main} reparam.}} = 1,
\end{equation}

In other words, $\pi(y\mid x)$ forms a proper probability distribution, meaning all probabilities are non-negative and their total sums to one. Examining Eq.~\ref{eq:op_policy}, we find that Eq.~\ref{eq:lag_p} corresponds to the partition function of the optimal policy that is determined by the reward function $r(x, y)$. The central insight underpinning the DPO algorithm is its ability to apply targeted constraints to the otherwise underdetermined Plackett-Luce (and notably Bradley-Terry) preference models. By doing so, it preserves the set of reward models that can be represented, while ensuring that the optimal policy described in Eq. \ref{eq:op_policy} remains analytically manageable for all input prompts $x$.

\subsection{Instability of Actor-Critic Algorithms} Our framework also enables an analysis of potential instabilities encountered by typical actor-critic methods in RLHF, such as PPO. Following the RLHF paradigm, our focus is on the RL fine-tuning phase described in Section \ref{section:prelims}. There exists a link to the control as inference framework \cite{levine2018reinforcement} in the context of the constrained RL problem in \ref{eq:RL}. Given a parameterized policy $\pi_{\theta}(y\mid x)$, the aim is to minimize $\mathbb{D}_{\text{KL}}[\pi_{\theta}(y|x) \mid \mid \pi^*(y\mid x)]$, where $\pi^*$ denotes the optimal policy derived from Eq. \ref{eq:optimum_model} and governed by the reward function $r_{\phi}(y, x)$. After some manipulation, the resulting objective can be written as:

\begin{equation}\label{eq:AC}
    \max_{\pi_{\theta}}\mathbb{E}_{\pi_{\theta}(y\mid x)}\bigg[\underbrace{r_{\phi}(x, y) -\beta\log\sum_{y}\pi_\text{ref}(y\mid x)\exp\left(\frac{1}{\beta}r_{\phi}(x, y)\right)}_{f(r_{\phi}, \pi_\text{ref}, \beta)} - \underbrace{\beta\log\frac{\pi_{\theta}(y\mid x)}{\pi_\text{ref}(y\mid x)}}_{\text{KL}}\bigg]
\end{equation}

The objective addressed here aligns with that explored in previous literature 
\citep{ziegler2020finetuning, stiennon2022learning, bai2022training, ouyang2022training}, employing the DPO-equivalent reward within the family of reward functions $r_{\phi}$. Within this framework, the normalization factor in $f(r_{\phi}, \pi_\text{ref}, \beta)$ can be viewed as representing the soft value function corresponding to the reference policy $\pi_\text{ref}$. Although this component leaves the optimal solution unchanged, omitting it can result in a policy gradient with substantial variance, which hampers stable learning. A possible remedy involves approximating the normalization factor with a learned value function, though this can introduce additional optimization challenges. Alternatively, previous studies have utilized a human completion baseline to normalize rewards—amounting to a one-sample Monte Carlo approximation of the normalization term. In contrast, the DPO reparameterization delivers a reward function that inherently avoids the need for baseline estimators.

\section{Experiments}
Here, we present empirical studies to assess the efficacy of DPO in training policies directly from preference data. We first investigate DPO’s sample efficiency and its capability to balance reward maximization against KL-regularization with the reference policy in a controlled text-generation task, comparing its performance to other prevalent preference optimization methods such as PPO. Subsequently, we extend our evaluation to larger models and more challenging RLHF tasks, including summarization and conversational modeling. Across these scenarios, we observe that DPO typically matches or surpasses the performance of robust baselines, such as RLHF with PPO and the selection of the highest-reward trajectory among $N$ candidates, requiring minimal hyperparameter adjustment. Details of our experimental protocol precede the presentation of results, with supplementary information provided in Appendix~\ref{app:exp_details}.

\textbf{Tasks.} Our study examines three distinct open-ended text generation scenarios. In all cases, algorithms aim to learn a policy from a dataset of preference annotations $\mathcal{D}=\bigl\{x^{(i)}, y_w^{(i)}, y_l^{(i)}\bigr\}_{i=1}^N$. For the \textbf{controlled sentiment generation} task, $x$ is an initial fragment of a movie review extracted from the IMDb dataset \cite{maas-EtAl:2011:ACL-HLT2011}, and the objective is for the policy to generate continuations $y$ that express positive sentiment. To enable a systematic evaluation, preference pairs between generated outputs are created automatically by leveraging a pre-trained sentiment classifier; a pair is selected such that $p(\text{positive}\mid x,y_w)>p(\text{positive}\mid x,y_l)$. The SFT baseline is established by fine-tuning GPT-2-large until convergence using movie reviews from the IMDb training split (with further implementation specifics provided in App~\ref{app:sentiment_details}). In the \textbf{summarization} setting, $x$ is a Reddit forum post and the model is tasked with producing a summary $y$ capturing the main ideas. In line with previous research, we utilize the Reddit TL;DR summarization dataset \citep{volske-etal-2017-tl} and incorporate human preference data collected by \citeauthor{stiennon2022learning}. The SFT model here is fine-tuned on human-authored summaries of forum posts\footnote{\url{https://huggingface.co/CarperAI/openai_summarize_tldr_sft}}, and reinforcement learning from human feedback (RLHF) is implemented using the TRLX framework \citep{leandro_von_werra_2023_7790115}. Notably, the human preference annotations come from samples produced by a different, but comparably fine-tuned, SFT model. Lastly, for the \textbf{single-turn dialogue} task, $x$ corresponds to a human prompt—ranging from scientific inquiries to personal advice requests. The policy is responsible for generating a response $y$ that is informative and engaging; for this, we draw upon the Anthropic Helpful and Harmless dialogue dataset \citep{bai2022training}, which comprises 170k conversations between humans and an automated assistant. Each dialogue concludes with a pair of model-generated responses and a human-annotated preference indicating which is superior. Because no SFT model pre-trained on this domain is available, we create one by fine-tuning a generic language model exclusively on those responses preferred by humans.

\textbf{Evaluation.} Our experimental design employs two primary evaluation methodologies. To assess the capacity of each algorithm for constrained reward maximization, we utilize a controlled sentiment generation scenario where performance is measured by charting the trade-off between realized reward and KL-divergence from a reference policy. This is feasible in this context due to the availability of an explicit reward function (a sentiment classifier). In contrast, such access to ground-truth reward is unavailable in real-world settings. Consequently, in tasks such as summarization and single-turn dialogue, we rely on the \textit{win rate} of each algorithm relative to a baseline policy, leveraging GPT-4 as a stand-in for human judges when evaluating summary quality or dialogue helpfulness. For summarization, the reference summary in the test set serves as the baseline, whereas in dialogue the baseline corresponds to the human-preferred response from the test set. While literature suggests that language models may surpass standard automatic metrics in evaluator efficacy \citep{Chen2023ExploringTU}, we substantiate our use of GPT-4 with a complementary human evaluation presented in Sec.~\ref{sec:human-judgments}. Our findings indicate a high degree of concordance between GPT-4 and human raters, often with human-GPT-4 agreement matching or exceeding the agreement between individual annotators. 

\textbf{Methods.} Beyond DPO, our analysis encompasses a variety of established methods for aligning language models with human preferences. We incorporate zero-shot prompting using \textbf{GPT-J} \citep{gpt-j} for summarization, and two-shot prompting with \textbf{Pythia-2.8B} \citep{biderman2023pythia} for dialogue tasks. Additionally, we include the \textbf{SFT} model as well as \textbf{Preferred-FT}, which is trained via supervised fine-tuning on the selected output $y_w$—sourced from the SFT model in sentiment control and summarization, or from a generic language model in single-turn dialogue. We also evaluate the \textbf{Unlikelihood}~\citep{welleck2019neural} approach, which explicitly increases the model's probability on $y_w$ and penalizes the likelihood of $y_l$, optionally scaling the unlikelihood loss with a factor $\alpha\in[0,1]$. Our benchmarking further includes \textbf{PPO} \citep{schulman2017proximal} with a reward function trained from preference data, and \textbf{PPO-GT}, an oracle variant that employs the ground-truth reward available in the controlled sentiment setting. For sentiment experiments, two PPO-GT implementations are tested: a standard out-of-the-box configuration \cite{leandro_von_werra_2023_7790115} and an adapted version incorporating reward normalization and refined hyperparameters for enhanced outcomes; these adjustments are likewise applied in `standard' PPO runs with learned rewards. As an additional comparator, we use the \textbf{Best of $N$} baseline, which samples $N$ candidate responses from the SFT model (or Preferred-FT in dialogue) and selects the response that maximizes a reward model trained on preference data. This approach, while often achieving strong performance, separates reward model quality from PPO optimization but is computationally demanding, as it necessitates generating $N$ samples for every input during testing.

\subsection{How well can DPO optimize the RLHF objective?}

The typical RLHF framework employs a KL-constrained reward maximization objective, aiming to maximize rewards while simultaneously limiting how much the learned policy diverges from a designated reference policy. Thus, in comparing different methods, it is critical to consider both the reward attained and the magnitude of KL divergence; a marginal increase in reward is not always preferable if it comes at the cost of a significantly greater KL value. In Figure~\ref{fig:frontier-tldr-main}, the reward-KL tradeoff frontiers for several methods are depicted in the sentiment analysis context. For each method, we conduct several independent training sessions, systematically varying the policy conservativeness hyperparameter (using target KL $\in\{3,6,9,12\}$ for PPO, $\beta \in \{0.05,0.1,1,5\}$ for DPO, $\alpha\in\{0.05,0.1,0.5,1\}$ for unlikelihood, and random seeds for preferred-FT). This experimental protocol leads to a total of 22 individual runs. At intervals of 100 training steps—up to convergence—we assess each resulting policy on a suite of held-out prompts, calculating both the mean reward using the true reward function and the average sequence-level KL divergence\footnote{Specifically, this is computed as the aggregate over per-timestep KL-divergences.} relative to the reference policy $\text{KL}\left(\pi\mid \mid \pi_\text{ref}\right)$. The empirical results demonstrate that DPO consistently achieves the most favorable frontier, obtaining superior rewards while maintaining comparatively low KL divergence. This is particularly striking for two main reasons. First, although DPO and PPO both target the same underlying objective, DPO’s frontier clearly outperforms that of PPO, strictly dominating its reward/KL tradeoffs. Second, DPO surpasses even PPO when PPO is supplied with ground truth rewards (PPO-GT), establishing its efficacy under identical conditions.

\subsection{Can DPO scale to real preference datasets?}
\label{sec:dpo-real-datasets}
We proceed to assess the effectiveness of DPO when fine-tuned on tasks involving summarization and single-turn dialogue. It is important to note that commonly used automatic metrics for summarization, like ROUGE, may not align well with human judgments~\citep{stiennon2022learning}. Previous studies have demonstrated that using PPO to align language models with human preferences can yield higher-quality summaries. To compare various methods, we generate model outputs from the test split of the TL;DR summarization dataset and calculate the mean win rate by pitting these outputs against the reference completions within the test set. Outputs for every approach are sampled using temperature values spanning from 0.0 to 1.0; the corresponding win rates are illustrated in Figure~\ref{fig:frontier-tldr-main} (right). DPO, PPO, and Preferred-FT are all adapted from the same GPT-J SFT model\footnote{\url{https://huggingface.co/CarperAI/openai_summarize_tldr_sft}}. Our analysis shows that DPO secures a win rate near 61\% at a temperature of 0.0, surpassing PPO, which reaches approximately 57\% at its best-performing sampling temperature. Additionally, DPO achieves a higher peak win rate in comparison to the best of $N$ baseline. It is worth mentioning that the $\beta$ hyperparameter for DPO was not extensively optimized, indicating that these findings may underrepresent DPO's capabilities. Furthermore, DPO maintains consistent performance across temperature variations, unlike PPO, whose effectiveness can diminish and revert to baseline GPT-J behavior at higher temperatures. Preferred-FT yields only marginal gains relative to the SFT model. In Section~\ref{sec:human-judgments}, we also present direct human preference comparisons between DPO and PPO, finding that at a temperature of 0.25, DPO samples were selected over PPO samples 58\% of the time.

For single-turn dialogue tasks, our evaluation focuses on the subset of the Anthropic HH dataset \citep{bai2022training} corresponding to instances featuring a single exchange between human and assistant. For GPT-4-based assessments, the analysis employs preferred completions in the test set as references when calculating the win rates of different training methods. Since an established SFT model for this scenario does not exist, we initialize with a pre-trained Pythia-2.8B and fine-tune it using Preferred-FT, training a reference model on the curated completions to ensure that generated responses remain in-distribution, subsequently applying DPO for further optimization. For benchmarking, we include both the top result among 128 Preferred-FT-generated completions (with Appendix Figure~\ref{fig:best-of-n} showing that performance gains plateau at $N=128$) and a 2-shot prompted version of the original Pythia-2.8B model. Across the evaluated temperatures, DPO achieves outcomes equal to or exceeding the strongest results among these baselines. Additionally, we investigate an RLHF model fine-tuned with PPO on the Anthropic HH dataset\footnote{\url{https://huggingface.co/reciprocate/ppo_hh_pythia-6B}}, sourced from a reputable implementation\footnote{\url{https://github.com/CarperAI/trlx/tree/main/examples/hh}}; however, we are unable to identify a prompt or temperature yielding improved performance over the base Pythia-2.8B. Based on observations from both the TL;DR setting and the shared objective between PPO and Best of 128, we use the latter as an approximate stand-in for PPO performance. In sum, DPO distinguishes itself as the sole efficient approach to meaningfully outperform preferred completions in the Anthropic HH test context and attains results comparable to or surpassing the computationally intensive Best of 128 baseline. As depicted in Figure~\ref{fig:dialogue-main}, DPO also reaches its optimal performance in relatively few training iterations.

\subsection{Generalization to a new input distribution}

To further investigate how PPO and DPO handle distribution shifts, we assess the performance of the policies trained on Reddit TL;DR summarization by testing them on a different dataset, specifically the news articles from the CNN/DailyMail test split \citep{nallapati-etal-2016-abstractive}, employing the best sampling temperatures identified on TL;DR (0 and 0.25). The outcomes, summarized in Table~\ref{tab:ood}, involve computing the GPT-4 win rate against the reference summaries using the same GPT-4 (C) evaluation prompt as in the Reddit TL;DR setup, with “forum post” replaced by “news article.” On this novel data distribution, DPO remains superior to the PPO policy, demonstrating a significant advantage. These findings provide preliminary support that DPO-trained policies exhibit comparable or better generalization capabilities than those trained with PPO, despite DPO not utilizing the extra unlabeled Reddit TL;DR data accessed by PPO.

\subsection{Validating GPT-4 judgments with human judgments}
\label{sec:human-judgments}
We conducted a human subject study to assess the validity of GPT-4’s evaluation, focusing on outputs from the TL;DR summarization experiments and employing two different prompting strategies. The \textbf{GPT-4 (S)} (simple) prompt straightforwardly requests an assessment of which summary more effectively covers the key information in the post. The \textbf{GPT-4 (C)} (concise) prompt adds an explicit preference for conciseness; we included this variant upon observing that, when using the \textbf{GPT-4 (S)} prompt, GPT-4 sometimes favors lengthier, more repetitive summaries compared to human evaluators. Full prompt texts are included in Appendix~\ref{app:prompts}. Our evaluation includes three experimental comparisons representing a range of summarization qualities: the highest performing method (DPO, temp. 0.25), the lowest performing (PPO, temp. 1.0), and a method with intermediate quality (SFT, temp. 0.25). Each is compared to the best-performing greedy PPO baseline. Analysis reveals that, across both prompt formulations, GPT-4's choices correspond with human preferences at rates similar to inter-human agreement, implying that GPT-4 serves as an adequate surrogate for human evaluators (multiple independent human assessments are gathered for the DPO and PPO-1 comparisons, though limited by available annotators). Across the board, win rates determined using the \textbf{GPT-4 (C)} prompt tend to align more closely with human evaluations; thus, this prompt is adopted for the central analyses in Section~\ref{sec:dpo-real-datasets}. Further methodological details—including the online interface provided to raters and a roster of study participants—are available in Appendix~\ref{app:human-study}.

\section{Discussion}
Preference-based learning offers a robust and scalable approach for developing competent and value-aligned language models. In this work, we presented DPO, a streamlined training framework that enables language models to learn directly from preference data, bypassing the need for reinforcement learning. Instead of adapting the preference learning problem to fit traditional RL paradigms and leveraging existing RL solutions, DPO establishes a direct correspondence between the language model’s policy and reward functions. This alignment allows for straightforward training with a cross-entropy loss, avoiding the complexities and potential limitations of reinforcement learning while preserving expressivity. Notably, DPO requires minimal hyperparameter tuning and matches or exceeds the performance of prominent RLHF techniques, such as those utilizing PPO, thereby making it considerably easier to train language models from human preference data.

\textbf{Limitations \& Future Work.} Our findings also highlight several avenues for continued investigation. A key question is how well the DPO-trained policy generalizes to data outside of its training distribution compared to models trained with explicit reward functions. Preliminary evidence points to generalization abilities on par with PPO-based methods, but systematic and in-depth evaluations are warranted. Further, it remains to be seen whether methods such as self-labeling—using the DPO policy to annotate previously unlabeled prompts—can be leveraged as effectively as in existing paradigms. Another open question involves the dynamics of reward over-optimization within direct preference optimization and whether the modest decline in performance observed in Figure~\ref{fig:dialogue-main}-right exemplifies this phenomenon. While our analysis focuses on models containing up to 6B parameters, extending DPO to much larger, cutting-edge architectures is a promising research path. In terms of assessment, our experiments reveal that GPT-4’s win rate evaluations can be sensitive to the prompt wording; future research could address optimal strategies for eliciting reliable evaluations from automated judges. Beyond language models, DPO holds promise for a broad spectrum of applications, such as the training of generative models across diverse data modalities.